\section{The Market, Value of Money, Securities and Strategies}
\subsection{Introduction}
\subsubsection{Terminologies}
\begin{definition}[Equity]
  The most basic of financial instruments is the \textbf{equity}, or known as the stock
  or the share.
\end{definition}
The \textbf{stockholder} (or \textbf{shareholder}) holds the right of
participating in the control of the company and receives \textbf{dividends}.
\begin{definition}[Dividend]
  \textbf{Dividends} are lump sum payments that are paid out every quarter or
  every six months to the stockholder, amount depending on the profitability of
  the company.
\end{definition}
\begin{definition}[Commodities]
  \textbf{Commodities} are raw products such as precious metals, oil and food products.
  Prices of commodities are unpredictable but tend to show seasonal effects,
  and are typically traded by people with no need for the material themselves.
\end{definition}

\subsection{Time Value of Money}
Depending on the market (in a deflating or inflating economy), but in general,
money today is has greater value than the same amount in a years time. This is
because money can be invested to ``produce'' more money, for instance through
interest payments.
\begin{remark}
  We consider investing money in a bank a \textbf{risk-free} investment and
  other financial instruments relative to this ``standard''.
\end{remark}
\begin{definition}[Effective Interest Rate]
  The interest rate associated with the basic time unit (typically defined $1$
  year) is the \textbf{effective interest rate} $r$.
\end{definition}
\begin{definition}[Nominal Interest Rate]
  When the frequency of compounding is not equal to the basic time unit, we call
  the associated interest rate the \textbf{nominal interest rate} $i$, defined:
  \[i=r\times \frac{t_{i}}{t_{r}}.\] 
  where $r$ is the effective interest rate, $t_{i}$ is the nominal interest
  time period, $t_{r}$ is the basic time unit of the effective interest rate.
\end{definition}
Suppose a nominal interest rate of $i$, with $m$ interest payments during a
year (once every $t=\frac{1}{m}$ years). The effective interest rate of each
sub-period $t$ years is $\frac{i}{m}$, and the effective interest rate for a
year is
\[r={(1+\frac{i}{m})}^{m}-1.\]
\subsubsection{Continuous Compounding}
As the interest payments happen at increasingly frequent intervals, $m\to\infty$:
\[\lim_{m \to \infty} {(1+\frac{r}{m})}^{m}= e^{r}.\]
\begin{proof}
  Using the binomial theorem,
  \begin{align*}
    {\left( 1+\frac{r}{m} \right)} ^{m}=&\sum_{i=0}^{m} \binom{m}{i}{\left( \frac{r}{m} \right)}^{i}\\
    =&1+\binom{m}{1}\frac{r}{m}+\binom{m}{2}{\left( \frac{r}{m} \right) }^{2}+\cdots+{\left( \frac{r}{m} \right)}^{m}\\
    =&1+r+\frac{m-1}{m}\frac{r^{2}}{2\mathpunct{!}}+\cdots+{\left( \frac{r}{m} \right) }^{m}
  \end{align*}
  Taking the limit $m\to\infty$:
  \[\lim_{m \to \infty} {(1+\frac{r}{m})}^{m}=1+r+\frac{r^{2}}{2\mathpunct{!}}+\frac{r^{3}}{3\mathpunct{!}}+\cdots= e^{r}.\]
  This is known as continuous compounding.
  \qed%
\end{proof}
\begin{remark}
  Continously compounded interest rates are used unless stated otherwise. The
  value of $b$ compounded over $T$ years at interest rate $r$ is: \[b\times
  e^{rT}.\]
\end{remark}
\subsubsection{Present Value}
\begin{definition}[Present Value]
  We define a mapping $P:\mathbb{R}\mapsto \mathbb{R}$ with initial condition
  $P(0)=k$ and interest rate $r$, such that:
  \[P(T)=ke^{rT}.\] 
  Due to the time value of money, cash flows in the future have smaller value
  that that in the present. We can then relate cash flows in the future (at
  time $T$) to their \textbf{present value} by:
  \[P(t)=P(T)e^{-r(T-t)},\quad t\in \left[ 0,T \right].\]
  This operation is known as \textbf{discounting}.
\end{definition}
\subsection{Fixed-Income Securities}
When a corporation wishes to borrow money from the public on a long-term basis,
it does so by issuing and selling debt securities known as \textbf{bonds}.
\begin{definition}[Bond]
  A \textbf{bond} is typically an interest-only loan, where interest is paid
  every period and principal is repaid at the end of the loan. It is defined by
  the \textit{face value/par value/nominal} (the lump sum), the
  \textit{coupons} and the \textit{maturity}.
\end{definition}
\begin{definition}[Coupons]
  The interest payments made on a bond is known as \textbf{coupons}. We can
  derive the \textit{coupon rate} from the \textit{annual coupon} $C$ and the
  face value $F$ of a bond $\frac{C}{F}$.
\end{definition}
The PV of a bond consists the PV of the face value and the annuity:
\begin{align*}
  PV_{\text{Bond}}=&PV_{\text{Face}}+PV_{\text{Coupon Cash Flow}}\\
  =&\frac{Face}{{(1+r)}^{t}}+\left(\frac{Annuity}{r}\right)\left( 1-\frac{1}{{(1+r)}^{t}} \right) 
\end{align*}
where $r$ is the interest rate, and $t$ the time to the bond's maturity.

\begin{definition}[Zero Coupon Bond]
  A \textbf{zero coupon bond} is a bond that pays no coupons at all. It must be
  offered at a price much lower than its face value.
\end{definition}

\subsection{Market Assumptions}
\begin{definition}[Arbitrage]
  Arbitrage refers to the possibility to make money without any risk.
\end{definition}
In a real market involving many participants, each motivated by making a
profit, any arbitrage opportunities that arises would vanish very quickly. We
make some market assumptions to ensure that there will only be a unique fair
price for a given bond.

\subsubsection{Terminologies}
\begin{definition}[Spot Price]
  We denote the price of any given stock or commodity at time $t\in
  \left.\left[0,\infty\right. \right)$ by $S(t)$. For present time (\textbf{spot time})
  $t=t_0$, $S(t_0)$ is the \textbf{spot price}.
\end{definition}
\begin{definition}[Price Process]
  For some $T\in \left( 0,\infty \right) $, the \textbf{price process} denoted ${\left(
  S(t) \right)}_{t\in [0,T]}=\left\{ S(t):t\in \left[ 0,T \right] \right\}$ for
  some function of randomness and time $S:\Omega\times \left[ 0,T
  \right]\mapsto [0,\infty)$. The price process is a random variable evolving
  through time, or a stochastic process.
\end{definition}
\begin{definition}[Underlyings]
  For contracts or derivatives written in relation to a certain amount of
  stocks or commodities, the particular stocks or commodities involved are
  referred to as the \textbf{underlyings}.
\end{definition}
\begin{definition}[Short-Selling]
  \textbf{Short-selling} involves selling an unowned or borrowed asset
  with the intention of buying it back later.
\end{definition}
\begin{definition}[Hedging]
  A strategy intended to reduce or minimise risk is known as \textbf{hedging}.
\end{definition}
\begin{definition}[Speculation]
  Taking a position in the market where the investor speculates that the price
  of the underlying asset will go up or down is known as \textbf{speculation}.
\end{definition}
\begin{definition}[Derivatives]
  \textbf{Derivatives} are financial contracts whose value depends from the
  values of some underlying asset(s).
\end{definition}

\subsubsection{Market Assumptions}
\begin{itemize}
  \item There are no transaction costs, e.g.\ no taxes, no bid-ask price, no
    brokerage fees. We call this market a \textbf{frictionless market}.
  \item The same risk-free rate of interest applies for borrowing and lending
    money (\textbf{market efficiency}) and the bank never goes bankrupt.
  \item There is an `absence' of arbitrage, as market participants take
    advantage of opportunities as they occur.
  \item It is always possible to buy/sell something in any amount at its fair
    price.
\end{itemize}

\subsection{Forward Contract}
\begin{definition}[Forward Contracts]
  \textbf{Forward Contracts} are agreements between two parties to buy or sell
  an asset $S(T)$ at a specified future time $T$ at a specified delivery price
  $K$, known as the maturity or delivery date.  Participants are obligated to
  buy or sell the asset at maturity.
\end{definition}
One of the parties assumes a long position to purchase the underlying asset at
time $T$, and the other party assumes a short position to sell the underlying
asset.
\subsubsection{Pay-off}
Pay-off per unit of asset from the long position in a forward contract: $S(T)-K$.\\
Pay-off per unit of asset from the short position in a forward contract: $K-S(T)$.
\begin{figure}[h!]
  \centering
  \includegraphics[width=0.8\textwidth]{forward-payoff}
\end{figure}\\

\centering
\begin{tabular}{c c c}
\toprule
Holding & Value at time $t$ & Value at maturity $T$\\
\midrule
Forward & $+0$ & $S(T)-K$\\
Stock & $-S(t)$ & $-S(T)$\\
Bank & $+S(T)$ & $S(t)e^{r(T-t)}$\\
\midrule
Total & $0$ & $S(t)e^{r(T-t)}-K$\\
\bottomrule
\end{tabular}

\raggedright%
\begin{prop}[Fair Delivery Price of a Forward Contract]
  The arbitrage free delivery price of a forward contract is a map $K: \left[
  0,T \right]\times [0,\infty)\mapsto [0, \infty)$ such that:
  \[K(t,T)=S(t)e^{r(T-t)}.\] 
\end{prop}
\begin{prop}[Convergence to Spot Price]
  The fair no-arbitrage delivery price of a forward contract converges to the spot price
  of the underlying asset as we approach maturity. That is,
  \[\lim_{t \to T} K(t,T)=S(T).\]
\end{prop}

\subsection{Options Contracts}
Options are similar to forward contracts, but without their obligatory aspect.
Additionally, options are sold for a certain specified initial amount, known as
the \textit{premium}.  The holder of a option has the decision whether to
execute the contract or not.  There are two types of options contracts:
\begin{itemize}
  \item A \textbf{call option} gives the holder the right to buy the underlying
    asset $S$ at a certain future \textit{expiration date} $T$ for a certain
    \textit{strike price} $K$.
  \item A \textbf{put option} gives the holder the right to sell the underlying
    asset $S$ at a certain future \textit{expiration date} $T$ for a certain
    \textit{strike price} $K$.
\end{itemize}
\begin{remark}
  There are two subtypes of the options contracts:
  \begin{itemize}
    \item \textbf{European options} can be exercised only on the expiration date itself.
    \item \textbf{American options} can be exercised at any time up to the expiration date.
  \end{itemize}
  Due to the additionally flexibility, premiums of the same parameters will be
  higher in American options than European options.
\end{remark}
\subsubsection{Pay-off}
Suppose an investor purchases a European call option on a certain stock with a
strike price $K=10$ and an expiration datae in three months.
At maturity, there are two possibilities for $S(T)$:
 \begin{enumerate}
  \item $S(T)<K$:
    The investor will choose not to exercise the option, as they can purchase
    the stock in the market for a better price. There is no profit to be made
    from the option and the investor loses money by paying for the premium.
  \item $S(T)>K$:
    The investor will exercise the option and will make a profit by selling them
    immediately at the spot price.
\end{enumerate} 
\begin{figure}[h!]
  \centering
  \includegraphics[width=0.8\textwidth]{option-long-payoff}
\end{figure}
The pay-off of the long European call is given by $\max \left\{ S(T)-K,0
\right\}$ and we define the \textbf{call function} the map
$f:\mathbb{R}\mapsto \mathbb{R}$:
\[f(x)=\max \left\{ x-K,0 \right\}={(x-K)}^{+}.\]
Similarly, the pay-off of the long European put is given by $\max \left\{
K-S(T),0 \right\}$ and we define the \textbf{put function} the map
$f:\mathbb{R}\mapsto \mathbb{R}$:
\[f(x)=\max \left\{ K-x,0 \right\}={(K-x)}^{+}.\]

\centering
\begin{tabular}{c c c}
  \toprule
  Contract & Position & Pay-off \\
  \midrule
  Call & Long & $\max \left\{ S(T)-K,0 \right\}$ \\
  Call & Short & $-\max \left\{ S(T)-K,0 \right\}$\\
  \midrule
  Put & Long & $\max \left\{ K-S(T) \right\}$\\
  Put & Short & $-\max \left\{ K-S(T),0 \right\}$ \\
  \bottomrule
\end{tabular}

\raggedright

\subsubsection{Replicating Portfolio}
\begin{definition}[Replicating Portfolios]
  We say that a portfolio is \textbf{replicating} with relation to a given
  contract if at maturity time, the future value of the replciating portfolio
  is the same as the future value of the contract.
\end{definition}
\begin{theorem}[Law of One Price]
  In an arbitrage free market, two portfolios with identical future values must
  have identical present values.
\end{theorem}
We see that replicating portfolios (by definition) have the same value at the
present and in the future.

\subsubsection{Put-Call Parity}
\begin{theorem}[Put-Call Parity]
  Consider a call and put option written on a stock $S$ with maturity $T$ and
  strike price $K$. Let ${\left( c_t \right)}_{t\in \left[ 0,T \right]}$ and
  ${\left( p_t \right)}_{t\in \left[ 0,T \right]}$ denote the fair prices of a
  call and put option at time  $t\in \left[ 0,T \right]$. Then for any $t\in \left[ 0,T \right]$,
  \[c_{t}+Ke^{-r(T-t)}=p_{t}+S(t).\] 
\end{theorem}
\begin{proof}
  Consider two portfolios:\\
  \textit{Portfolio A}: A cash amount of £$Ke^{-rT}$ and a European call option
  written on a non-dividend-paying stock.\\
  \textit{Portfolio B}: A European put option written on a non-dividend-paying
  stock and one share.
  At time $t=T$, we have:
  \[V_A=\max \left\{ S(T),K \right\}\qquad V_B=\max \left\{ S(T),K \right\}.\]
  They are replicating portfolios and by the law of one price, the values
  of the portfolios at time $t=0$ should be equal. That is:
  \[c_{0}+Ke^{-rT}=p_{0}+S(0).\]
  We obtain the put-call parity at time $t=0$ which applies $\forall t\in [0,T]$ by finding the
  values of the portfolios at time $t$.
  \qed%
\end{proof}

\subsubsection{Synthetic Securities}
Using the put-call parity, we may create synthetic versions of other
securities, of which replicates or simulate a security by combining the
features of other assets.
\begin{example}[Synthetic Stock]
  A long position in a underlying stock is equivalent to a long position in a
  call option, a long position in a risk-free bond and a short position in a
  put option.
  \[S(0)=c_0+Ke^{-rT}-p_0.\] 
\end{example}
\begin{example}[Synthetic Put Option]
  A long position in a put option is equivalent to a long position in a call
  option, a long position in a risk-free bond and a short position in the
  underlying stock.
  \[p_0=c_0+Ke^{-rT}-S(0).\] 
\end{example}

\subsection{Trading Strategies}
With these financial objects, we may derive further common strategies driven by
speculation motives or desires to hedge, of which can
\begin{enumerate}
  \item involve a single option and a stock
  \item involve two or more options of the same type (\textbf{spreads})
  \item involve both calls and puts on the same stock (\textbf{combinations})
\end{enumerate}
\begin{remark}
  We ignore the time value of money in calculating the profit functions, that
  is, $r=0$. Profits are calculated simply by subtracting the initial cost from
  the final payoff.
\end{remark}
\subsubsection{Covered Calls, Protective Puts}
These strategies seek to reduce the risk from the short position of a call
option, or the risk of a sharp decrease in the value of a stock. We hedge the
unlimited/large downside risk of large price changes, by holding a long
position in the underlying stock, or the long position in a put option.
\begin{figure}[h!]
  \centering
  \includegraphics[width=0.8\textwidth]{covered-protective}
\end{figure}
An investor may also short these financial objects, of which the profit
diagrams is the negative of those above.
\subsubsection{Spreads}
A spread is a portfolio consisting two or more options of the same type, with
different positions.
\begin{itemize}
  \item \textbf{Bull Spread}: long one call with strike price $K_l$, and short
    one call with higher strike price $K_s>K_l$ 
  \item \textbf{Bear Spread}: long one call with strike price $K_l$, and short
    one call with lower strike price $K_s<K_l$ 
  \item \textbf{Butterfly Spread}: long two calls with strike prices
    $K_{l1},K_{l2}$ and short two calls with strike price $K_s = \frac{K_{l1} +
    K_{l2}}{2}$
\end{itemize}
\subsubsection{Combinations}
A combination is a portfolio consisting both call and put options with the same
position on the same stock. The strikes and maturities of all the options
included are the same.
\begin{itemize}
  \item \textbf{Straddle}: long one call and one put
  \item \textbf{Strips}: long one call and two puts
  \item \textbf{Straps}: long one put and two calls
  \item \textbf{Strangles}: long a put and a call
\end{itemize}
